\chapter[Historical Notes]{Historical Notes}\label{ch:sec:historical-notes}

This chapter collects historical remarks that contextualise the
theory of weighted incidence structures within the broader
development of cryptography, information theory, and geometry.
Each section corresponds to a part of the book.

\section{The Origins of Universal Optimality}\label{ch:sec:historical-uo}

The concept of universal optimality was first identified by
Shannon~\cite{shannon1949} under the name of \emph{perfect secrecy},
characterised by the condition $\Pr[M\mid C]=\Pr[M]$.
Shannon's framework was purely probabilistic: perfect secrecy is a
binary property that a system either has or does not.

The transition from a binary to a continuous measure of secrecy
began with Wyner~\cite{wyner1975wiretap}, who introduced the
wiretap channel and showed that secrecy could be quantified as
the equivocation rate.  Maurer~\cite{maurer1993secret} extended
this to the secret-key agreement model, and Csisz\'ar~\cite{csiszar1996almost}
developed the information-theoretic approach to universal
optimality in the context of source coding.

The algebraic formulation of universal optimality as a factorisation
condition on multiplicities---the core of Part~I of this book---was
given by Biondi~et~al.~\cite{biondi2018apollonian}, who also
introduced the term \emph{universal optimality} and proved the
representation theorem for deterministic encoders.

\section{Weighted Incidence Structures}\label{ch:sec:historical-wis}

Weighted incidence structures generalise the combinatorial notion
of a bipartite graph with vertex weights.  The factorisation
$n_{mc}=\beta(c)/w(m)$ appears implicitly in the study of
Apollonian cell encoders~\cite{biondi2018apollonian}, where the
weights $w,\beta$ arise from integer partitions of the key space.

The gauge symmetry $w\mapsto\lambda w,\;\beta\mapsto\lambda\beta$
was recognised by F"{u}ster-Sabater and colleagues in the context
of side-channel analysis, where different normalisations of the
power trace produce equivalent leakage models.

The connection between weighted incidence structures and graph
Laplacians---developed in Part~II---builds on the spectral graph
theory of Chung~\cite{chung1997spectral} and the theory of
Schur complements for bipartite graphs, which originated in
the study of electrical networks~\cite{vanvalkenburg1974network}.

\section{Posterior Manifold and Information Geometry}\label{ch:sec:historical-posterior}

The posterior manifold constructed in Part~III is a Riemannian
submanifold of the probability simplex.  The use of logarithmic
coordinates for the simplex dates back to Aitchison~\cite{aitchison1982statistical}
in compositional data analysis, where log-ratios are the natural
coordinates for proportions.

The pullback metric through the logarithmic map was studied
by Amari~\cite{amari2016information} as the \emph{self-information
metric} or \emph{Fisher metric} for exponential families with
the log-partition function as potential.  However, our metric
differs from the Fisher metric because it is pulled back through
the logarithmic map rather than the square-root map, reflecting
the underlying additive structure of the WIS factorisation rather
than the estimation-theoretic origin of the Fisher information.

\section{Universal Bounds and Bregman Divergences}\label{ch:sec:historical-bounds}

The Lipschitz and quadratic bounds of Part~IV belong to the
tradition of \emph{geometric inequalities} in information theory,
exemplified by Pinsker's inequality $\operatorname{KL}(P\|Q)\ge\frac12\|P-Q\|_{\mathrm{TV}}^2$
and the reverse inequality of Csisz\'ar~\cite{csiszar1996almost}.
Our bounds are sharper because they are derived from the intrinsic
geometry of the posterior manifold rather than from general
inequalities between divergences.

The Average Hessian Operator was introduced by
Bregman~\cite{bregman1967relaxation} in the context of convex
optimisation, where it characterises the second-order behaviour
of Bregman divergences.  Its application to the posterior manifold
is new and provides a uniform framework for bounding the
variation of information-theoretic quantities.

\section{Variational Theory and Spectral Methods}\label{ch:sec:historical-variational}

The Three-Level Theorem (Part~V) belongs to the variational theory
of quadratic energies on simplices, whose roots lie in the
calculus of variations and the theory of Chebyshev systems.
The collapse to at most three distinct values is reminiscent of
the classical three-distance theorem in Diophantine approximation
and the three-gap theorem in dynamical systems.

The connection between the posterior Laplacian and spectral
graph theory establishes a bridge between the geometry of
encryption systems and the spectral theory of random walks on
graphs~\cite{chung1997spectral}.  The distribution-dependent
nature of our Laplacian---it varies with the posterior---has
no direct analogue in classical spectral graph theory.

\section{Open Questions and Future Directions}\label{ch:sec:historical-open}

The theory developed in this book raises as many questions as it
answers.  The extension to continuous message spaces, the quantum
generalisation of weighted incidence structures, and the
computational complexity of constructing universally optimal
encoders are among the most pressing open problems.
The final chapter of the book (Chapter~\ref{ch:sec:open})
provides a detailed research programme organised by priority.